%% kapa: exam=1 round=28 year=2017 period=1 form=A subject=real_estate subject_ko=부동산학원론 seq=29 number=109
%% source: Q-Net 감정평가사 기출문제 (cst00309) artlSeq=5199603 file=extracted/1008/5199603/2017년 제28회 감정평가사 1차시험 문제지/1교시 A형.pdf
%% conversion: kordoc 4.24.51 -> tools/gen1.py (auto); review=proofread-sonnet; flags=tikz
\begin{question}{109}
A와 B도시 사이에 C마을이 있다. 레일리의 소매인력법칙을 적용할 경우, C마을에서 A도시와 B도시로 구매 활동에 유인되는 인구수는? (단, C마을 인구의 60\%만 A도시 또는 B도시에서 구매하고, 주어진 조건에 한함)
\begin{center}
\begin{tikzpicture}[x=0.12mm,y=-0.12mm,font=\small,line width=0.7pt]
\draw (70,107) ellipse (65 and 72);
\node[align=center] at (70,107) {A도시\\10만명};
\draw (135,107) -- (258,107);
\node at (195,136) {20Km};
\draw (258,62) rectangle (375,153);
\node[align=center] at (316,107) {C마을\\2만명};
\draw (375,107) -- (488,107);
\node at (432,136) {20Km};
\draw (599,105) ellipse (111 and 100);
\node[align=center] at (599,105) {B도시\\30만명};
\end{tikzpicture}
\end{center}
\begin{choices}
  \item A: 3천명, B: 9천명
  \item A: 4천명, B: 8천명
  \item A: 5천명, B: 7천명
  \item A: 5천5백명, B: 6천5백명
  \item A: 6천명, B: 6천명
\end{choices}
\end{question}
\answer{1}
